\documentclass[11pt]{article}
\usepackage[margin=1in]{geometry}
\usepackage[T1]{fontenc}
\usepackage[utf8]{inputenc}
\usepackage{lmodern,microtype,amsmath,amssymb,amsthm,mathtools}
\usepackage[colorlinks=true,linkcolor=blue,citecolor=blue,urlcolor=blue]{hyperref}
\hypersetup{pdftitle={Minimum Gaussian energy for parity requires exponential ReLU width},pdfauthor={ChatGPT-6 Astra; Kaarel Hanni; Dmitry Vaintrob}}
\newtheorem{theorem}{Theorem}
\newtheorem{lemma}[theorem]{Lemma}
\newtheorem{corollary}[theorem]{Corollary}
\newcommand{\R}{\mathbb R}
\newcommand{\E}{\mathbb E}
\newcommand{\relu}{\mathrm{ReLU}}
\newcommand{\chiS}{\chi}
\title{Minimum Gaussian energy for parity requires\protect\\exponential ReLU width}
\author{ChatGPT-6 Astra \and Kaarel H\"anni \and Dmitry Vaintrob\thanks{Authors are listed in alphabetical order.}}
\date{6 November 2026}
\begin{document}
\maketitle
\begin{abstract}
A Boolean function may have a small neural-network implementation while
its minimum-energy implementation requires much greater width.
We prove an exact example for parity and one-hidden-layer ReLU networks
with a constant input coordinate and independent isotropic Gaussian
weights in each layer. For an even parity of degree $r\ge4$, the smallest
Gaussian rate energy compatible with output margin $1/2$ is attained
if and only if the hidden width is at least $2^{r-1}$. Nevertheless the
same parity has a bounded-weight shallow teacher of width $O(r^2)$ and
a Boolean circuit of size $O(r)$.
The proof first determines the sharp correlation of a single ReLU neuron
with parity, then characterizes every equality case by its Fourier
coefficients. The result concerns exact energy minimization. It is not
an exponential sample-complexity or online-mistake lower bound.
\end{abstract}

\section{The question and the result}

For a Gaussian prior on network weights, the leading small-noise
probability of a prescribed function is determined by the least
quadratic weight energy that implements it. The width needed to
\emph{implement} a function need not be the width needed to minimize
that energy. This distinction matters when comparing strong-prior
Bayesian learners at different widths: adding neurons can change the
leading function probabilities even when the target was already
realizable.

We give an exact separation for parity. Let $d\ge r\ge2$, let $r$ be even,
and write
\[
 X=\{-1,1\}^d,\qquad \chi(x)=\prod_{i=1}^r x_i.
\]
The learner has one hidden layer and a linear readout:
\begin{equation}\label{eq:model}
 z(x)=\sum_{j=1}^N v_j\relu(b_j+w_j\cdot x),
 \qquad \relu(t)=\max\{t,0\}.
\end{equation}
The bias $b_j$ is the weight of an appended constant input coordinate.
For the independent Gaussian prior
\[
 (b_j,w_{j1},\ldots,w_{jd})_i\sim N(0,N\sigma^2/d),
 \qquad v_j\sim N(0,\sigma^2/N),
\]
the density is proportional to $\exp[-K_N/(2\sigma^2)]$, where
\begin{equation}\label{eq:energy}
 K_N=\frac dN\sum_{j=1}^N(b_j^2+\|w_j\|_2^2)
       +N\sum_{j=1}^Nv_j^2.
\end{equation}
All first-layer coordinates, including the constant coordinate, have
the same variance. Impose the full-cube margin constraints
\begin{equation}\label{eq:margin}
 \chi(x)z(x)\ge\frac12\qquad(x\in X).
\end{equation}
Set
\begin{equation}\label{eq:rho}
 \rho_r=\frac{\binom{r-2}{(r-2)/2}}{2^{r-1}\sqrt r},
 \qquad E_* =\frac{\sqrt d}{\rho_r}.
\end{equation}

\begin{theorem}\label{thm:main}
Every width and every network satisfying \eqref{eq:margin} obey
$K_N\ge E_*$. If $r\ge4$, equality is attainable exactly at the widths
\[
 N\ge2^{r-1}.
\]
For $r=2$, equality is attainable exactly when $N\ge4$.
At any smaller width where \eqref{eq:margin} is feasible, its minimum
energy is strictly greater than $E_*$.

The target has a bounded-fan-in Boolean circuit of size $O(r)$ and a
one-hidden-layer ReLU teacher of width $O(r^2)$, every weight at most
one in absolute value, and output margin one.
\end{theorem}

The strict-margin cell $\chi(x)z(x)>1/2$ has the same energy infimum as
\eqref{eq:margin}: multiply the readout by $1+\varepsilon$ and let
$\varepsilon\downarrow0$. Thus the theorem also identifies the width
at which that cell reaches its optimal leading Gaussian small-noise
rate. The claim is about attaining the \emph{exact} minimum. It gives no
positive lower bound, uniform in $r$, on the approximation gap below
the width threshold.

Prior work establishes $\Theta(r)$ variation-norm complexity for parity
and shows that multidirectional ReLU representations improve on ridge
representations~\cite{ardeshir}. Our proof uses the augmented Euclidean
norm, which charges the bias, and determines the sharp single-neuron
constant and the exact support size needed for equality.

\section{A sharp single-neuron inequality}

For $T\subseteq[d]$, put $\chi_T(x)=\prod_{i\in T}x_i$ and
$\widehat h(T)=\E[\chi_T(x)h(x)]$, with expectation uniform on the cube.

\begin{lemma}[Top Fourier coefficient of a monotone function]
\label{lem:monotone}
If $h:\{-1,1\}^m\to[0,1]$ is coordinatewise nondecreasing and $m\ge1$,
then
\[
 \left|\E\!\left[\left(\prod_{i=1}^m x_i\right)h(x)\right]\right|
 \le\frac{\binom{m-1}{\lfloor(m-1)/2\rfloor}}{2^m}.
\]
\end{lemma}
\begin{proof}
Average $h$ over coordinate permutations. The objective is unchanged,
and the resulting value $h_j$ at inputs with $j$ positive coordinates
is nondecreasing in $j$: a uniform $j$-subset can be coupled with a
uniform $(j+1)$-subset by adding a uniformly chosen absent coordinate.
Every sequence $0\le h_0\le\cdots\le h_m\le1$ is a convex combination
of the two constant sequences and the threshold sequences
$\mathbf1\{j\ge t\}$. For a threshold, the alternating sum is
\[
 \sum_{j=t}^m(-1)^j\binom mj=(-1)^t\binom{m-1}{t-1}
 \qquad(1\le t\le m).
\]
The largest absolute value is the central binomial coefficient in the
statement. The constant sequences have zero top coefficient.
\end{proof}

\begin{lemma}[Sharp parity correlation]\label{lem:correlation}
For even $r\ge2$ and every $(b,w)\in\R^{d+1}$,
\begin{equation}\label{eq:correlation}
 \left|\E[\chi(x)\relu(b+w\cdot x)]\right|
 \le \frac{\binom{r-2}{(r-2)/2}}{2^{r-1}}
       \min_{i\le r}|w_i|
 \le\rho_r\sqrt{b^2+\|w\|_2^2}.
\end{equation}
For unit augmented norm, equality in the last bound holds exactly at
\[
 b=0,\qquad w_i=s_i/\sqrt r\ (i\le r),\qquad
 w_i=0\ (i>r),\qquad s_i\in\{-1,1\}.
\]
\end{lemma}
\begin{proof}
Input sign changes make the relevant $w_i$ nonnegative and change at
most the sign of the correlation. Choose one of their smallest values,
$a=w_j$. Condition on the irrelevant inputs. Since the derivative of
ReLU is the positive-half-line indicator almost everywhere,
\begin{align*}
 \E[\chi\relu(b+w\cdot x)]
 =\frac12\int_{-a}^a
 \E\left[\left(\prod_{i\le r,\ i\ne j}x_i\right)
 \mathbf1\left\{b+t+\sum_{i\ne j}w_ix_i>0\right\}\right]dt.
\end{align*}
The inner function is monotone in the remaining $r-1$ relevant bits.
Lemma~\ref{lem:monotone} bounds its top coefficient by
$\binom{r-2}{(r-2)/2}/2^{r-1}$, proving the first inequality.
Next,
\[
 \min_{i\le r}|w_i|\le
 \frac{\|(w_1,\ldots,w_r)\|_2}{\sqrt r}
 \le\frac{\sqrt{b^2+\|w\|_2^2}}{\sqrt r}.
\]
Equality forces the stated directions. They attain equality: at equal
weights and bias zero, the inner threshold in the integral is the same
central Hamming threshold for almost every $t\in(-a,a)$. Its top
coefficient has constant sign and maximal magnitude.
\end{proof}

For completeness, an odd parity of degree $r\ge3$ has sharp constant
$\rho_{r+1}$ when the bias is charged in the augmented norm. Indeed
$\relu(t)=(t+|t|)/2$, the affine part has zero parity correlation, and
randomizing the sign of the bias converts the absolute-value
correlation into an even $(r+1)$-variable parity correlation. Equality
is attained when the bias and all relevant weights have magnitude
$1/\sqrt{r+1}$. We will only need the even case below.

\section{Energy equality determines the output and all neuron directions}

The margins imply $\E[\chi z]\ge1/2$. Apply
Lemma~\ref{lem:correlation} separately to each neuron and then AM--GM:
\begin{align}\label{eq:lower}
 \frac12&\le\E[\chi z]
 \le\rho_r\sum_j|v_j|\sqrt{b_j^2+\|w_j\|_2^2},\\
 K_N&\ge2\sqrt d\sum_j|v_j|\sqrt{b_j^2+\|w_j\|_2^2}
 \ge\frac{\sqrt d}{\rho_r}.\nonumber
\end{align}
This proves the universal lower bound.

If equality holds, every nonzero neuron attains the sharp correlation
bound, with its readout sign making that correlation positive. A neuron
that contributes zero must have all parameters zero. Also
$\E[\chi z]=1/2$ and the pointwise margins force
\begin{equation}\label{eq:exact-output}
 z(x)=\chi(x)/2\qquad(x\in X).
\end{equation}
Write
\[
 h_s(x)=\relu\left(\frac1{\sqrt r}\sum_{i=1}^rs_ix_i\right),
 \qquad \kappa_r=\widehat h_{\boldsymbol1}([r]),\qquad
 |\kappa_r|=\rho_r.
\]
Merging proportional neurons, every equality case has the form
\begin{equation}\label{eq:measure}
 z(x)=\operatorname{sign}(\kappa_r)V_*
       \sum_s p_s\chi(s)h_s(x),
 \qquad p_s\ge0,\quad\sum_sp_s=1,\quad
 V_* =\frac1{2\rho_r}.
\end{equation}
Each orientation with $p_s>0$ requires at least one nonzero neuron.

\section{Fourier coefficients force exponentially many orientations}

The equality-case networks ignore the other $d-r$ inputs. Throughout
this section, Fourier coefficients and orientation characters are indexed
by subsets of the $r$ relevant coordinates.

\begin{lemma}\label{lem:ridge}
Let $h=h_{\boldsymbol1}$. Its constant coefficient is positive, every
even-degree coefficient is nonzero, and its odd-degree coefficients
vanish except at degree one.
More explicitly, if $|T|=2t\ge2$, $n=r-2$, and $m=2t-2$, then
\begin{equation}\label{eq:equator}
 \widehat h(T)=
 \frac1{2\sqrt r}\frac{\binom n{n/2}}{2^n}
 (-1)^{m/2}\frac{\binom{n/2}{m/2}}{\binom nm}.
\end{equation}
\end{lemma}
\begin{proof}
The identity $h(x)-h(-x)=\sum_i x_i/\sqrt r$ gives the odd-degree
claims, including $\widehat h(\{i\})=1/(2\sqrt r)$.
For $i,j\in T$, average first over these two coordinates. The remaining
sum $R=\sum_{\ell\ne i,j}x_\ell$ is even, and direct evaluation gives
\[
 \E_{x_i,x_j}[x_ix_jh(x)]
 =\frac1{2\sqrt r}\mathbf1\{R=0\}.
\]
On a balanced $n$-bit input,
$\prod_{\ell=1}^n(1+u x_\ell)=(1-u^2)^{n/2}$.
Comparing coefficients and using permutation symmetry shows that the
conditional mean of any degree-$m$ character on that equator is
$(-1)^{m/2}\binom{n/2}{m/2}/\binom nm$. Multiplying by its probability
$2^{-n}\binom n{n/2}$ proves \eqref{eq:equator}.
\end{proof}

Taking Fourier coefficients in \eqref{eq:measure},
\begin{equation}\label{eq:fourier-output}
 \widehat z(T)=\operatorname{sign}(\kappa_r)V_*\widehat h(T)
      \E_p[\chi_{[r]\mathbin\triangle T}(s)].
\end{equation}
For every even $T\ne[r]$, the ridge coefficient is nonzero while
\eqref{eq:exact-output} makes the left side zero. Hence
\[
 \E_p\chi_U=0\qquad\text{for every nonempty even }U\subseteq[r].
\]
The antipodal symmetrization of $p$ has these even Fourier moments and
zero odd Fourier moments. Fourier inversion makes it uniform, so
\begin{equation}\label{eq:pair-mass}
 p_s+p_{-s}=2^{1-r}\qquad(s\in\{-1,1\}^r).
\end{equation}
Every antipodal pair therefore needs a nonzero neuron. This proves
$N\ge2^{r-1}$.

For $r=2$, the nonzero degree-one ridge coefficients in
\eqref{eq:fourier-output} additionally force both singleton moments of
$p$ to vanish. All its nonconstant Fourier moments vanish, so $p$ is
uniform on all four orientations. Thus $N\ge4$ in that case.

\section{Matching networks and a small bounded-weight teacher}

For even $r\ge4$, let $p$ be uniform on $s_1=+1$. Its only nonzero
Fourier moments are those of $\varnothing$ and $\{1\}$. The empty
moment gives the desired parity coefficient in
\eqref{eq:fourier-output}. The moment $\{1\}$ would give a coefficient
at $T=[r]\setminus\{1\}$; its degree is odd and at least three, so the
ridge coefficient vanishes. Therefore \eqref{eq:measure} is exactly
$\chi/2$ and uses $2^{r-1}$ atoms. For $r=2$, use all four orientations
uniformly.

Write the nonzero atom coefficients as $c_s$. At any width $N$ at least
their number, set
\[
 (b_s,w_s)=t_s(0,s/\sqrt r,0,\ldots,0),\qquad
 v_s=c_s/t_s,\qquad t_s^2=\frac{N|c_s|}{\sqrt d}.
\]
Positive homogeneity preserves the output, and the two quadratic costs
of each neuron agree. Thus the total energy is
$2\sqrt d\sum_s|c_s|=E_*$. Zero neurons fill any additional width.

At a smaller feasible width the closed margin problem attains its
minimum, because the constraints are closed and $K_N$ is coercive.
The equality characterization excludes $E_*$, so the minimum is
strictly larger. This completes the energy and width assertions.

Finally, parity has a small bounded-weight teacher. Let
$S(x)=\sum_{i=1}^r x_i$ and linearly interpolate the alternating values
$\pm1$ at $S=-r,-r+2,\ldots,r$. The resulting one-variable spline has
$O(r)$ hinges, initial slope of magnitude one, slope jumps of magnitude
at most two, and constant term $O(r)$. Its hinge expansion gives a
shallow ReLU representation of parity. Each affine form has input
coefficients at most one and bias at most $r$. Divide it by the largest
absolute coefficient and replicate the resulting readout coefficient
into $O(r)$ copies, each of magnitude at most one. The constant and
linear terms take $O(r)$ further neurons. All incoming and outgoing
weights are then bounded by one, the total width is $O(r^2)$, and the
output is exactly $\chi$, with margin one. An ordinary XOR chain gives
the $O(r)$ Boolean circuit.

\section{Interpretation}

The theorem exhibits an exact difference between realizability width
and minimum-energy width under an unchanged isotropic Gaussian prior.
It does not show that a Bayesian learner needs exponential width to
learn parity, nor that an approximation to the optimal energy requires
exponential width. In particular, no online input sequence or mistake
lower bound is asserted. The result identifies a concrete reason why
an energy calculation for a narrow network need not determine the
strong-prior posterior at larger widths.

\begin{thebibliography}{9}
\bibitem{ardeshir}
Navid Ardeshir, Daniel Hsu, and Clayton Sanford.
Intrinsic dimensionality and generalization properties of the
$\mathcal R$-norm inductive bias.
\emph{Proceedings of the 36th Conference on Learning Theory},
PMLR 195, 2023.
\url{https://proceedings.mlr.press/v195/ardeshir23a.html}.
\end{thebibliography}
\end{document}
